% GENERATED FILE. Edit canonical lessons, then run npm run generate:books.
% canonical-source-hash: fnv1a64:71ff32a86e41cdea
% canonical-lessons: TE-C314-mestri, TE-C314-plambaru, TE-C314-mekanikku, TE-C314-kapaladaru, TE-C314-palavadu, TE-R314-first-pass-words-from-chapters-310-311, TE-R314-second-pass-words-from-chapters-312-314

\chapter{Mason, Plumber, Mechanic, Watchman, Milkman}
\label{ch:te-mason-plumber-mechanic-watchman-milkman}
\hlchaptermodality{314}

\begin{chapteropening}
\textbf{By the end of this chapter:} \emph{I can say a mason, a plumber, a mechanic, a watchman and a milkman.}

Complete the last lesson of chapter 314: I can say a mason, a plumber, a mechanic, a watchman and a milkman.
\end{chapteropening}

\section[\texorpdfstring{\te{మేస్త్రి} (mēstri)}{మేస్త్రి (mēstri)}]{\te{మేస్త్రి} (mēstri) — a mason, a foreman}

\label{lesson:TE-C314-mestri}



\begin{quote}
Before the new one: say the Telugu for ink, then the Telugu for attendance.
\end{quote}

\subsection*{You'll want to know: \te{మేస్త్రి}}
\textbf{\te{మేస్త్రి}} — \emph{mēstri} — \textquotedblleft{}a mason, a foreman\textquotedblright{}.

\begin{grammarlens}[title={What you've built}]
One more word for this chapter.
\end{grammarlens}

\subsection*{Your turn}
\begin{itemize}
\raggedright
  \item \emph{Say it:} \emph{mēstri}
  \item \emph{Say it:} \emph{mēstri}, once more
  \item \emph{Say it:} say \emph{mēstri}
  \item \emph{Recall:} say the Telugu for ink, then the Telugu for attendance, then say \emph{mēstri} again
\end{itemize}

\begin{tcolorbox}[breakable,colback=teal!4,colframe=teal!35!black,title={Before you move on}]
What does \te{మేస్త్రి} mean? (\textquotedblleft{}A mason, a foreman\textquotedblright{}.) Say it once more.
\end{tcolorbox}
\section[\texorpdfstring{\te{ప్లంబరు} (plambaru)}{ప్లంబరు (plambaru)}]{\te{ప్లంబరు} (plambaru) — a plumber}

\label{lesson:TE-C314-plambaru}



\begin{quote}
Before the new one: say the Telugu for attendance, then the Telugu for a mason.
\end{quote}

\subsection*{You'll want to know: \te{ప్లంబరు}}
\textbf{\te{ప్లంబరు}} — \emph{plambaru} — \textquotedblleft{}a plumber\textquotedblright{}.

\begin{grammarlens}[title={What you've built}]
One more word for this chapter.
\end{grammarlens}

\subsection*{Your turn}
\begin{itemize}
\raggedright
  \item \emph{Say it:} \emph{plambaru}
  \item \emph{Say it:} \emph{plambaru}, once more
  \item \emph{Say it:} say \emph{plambaru}
  \item \emph{Recall:} say the Telugu for attendance, then the Telugu for a mason, then say \emph{plambaru} again
\end{itemize}

\begin{tcolorbox}[breakable,colback=teal!4,colframe=teal!35!black,title={Before you move on}]
What does \te{ప్లంబరు} mean? (\textquotedblleft{}A plumber\textquotedblright{}.) Say it once more.
\end{tcolorbox}
\section[\texorpdfstring{\te{మెకానిక్కు} (mekānikku)}{మెకానిక్కు (mekānikku)}]{\te{మెకానిక్కు} (mekānikku) — a mechanic}

\label{lesson:TE-C314-mekanikku}



\begin{quote}
Before the new one: say the Telugu for a mason, then the Telugu for a plumber.
\end{quote}

\subsection*{You'll want to know: \te{మెకానిక్కు}}
\textbf{\te{మెకానిక్కు}} — \emph{mekānikku} — \textquotedblleft{}a mechanic\textquotedblright{}.

\begin{grammarlens}[title={What you've built}]
One more word for this chapter.
\end{grammarlens}

\subsection*{Your turn}
\begin{itemize}
\raggedright
  \item \emph{Say it:} \emph{mekānikku}
  \item \emph{Say it:} \emph{mekānikku}, once more
  \item \emph{Say it:} say \emph{mekānikku}
  \item \emph{Recall:} say the Telugu for a mason, then the Telugu for a plumber, then say \emph{mekānikku} again
\end{itemize}

\begin{tcolorbox}[breakable,colback=teal!4,colframe=teal!35!black,title={Before you move on}]
What does \te{మెకానిక్కు} mean? (\textquotedblleft{}A mechanic\textquotedblright{}.) Say it once more.
\end{tcolorbox}
\section[\texorpdfstring{\te{కాపలాదారు} (kāpalādāru)}{కాపలాదారు (kāpalādāru)}]{\te{కాపలాదారు} (kāpalādāru) — a watchman, a guard}

\label{lesson:TE-C314-kapaladaru}



\begin{quote}
Before the new one: say the Telugu for a plumber, then the Telugu for a mechanic.
\end{quote}

\subsection*{You'll want to know: \te{కాపలాదారు}}
\textbf{\te{కాపలాదారు}} — \emph{kāpalādāru} — \textquotedblleft{}a watchman, a guard\textquotedblright{}.

\begin{grammarlens}[title={What you've built}]
One more word for this chapter.
\end{grammarlens}

\subsection*{Your turn}
\begin{itemize}
\raggedright
  \item \emph{Say it:} \emph{kāpalādāru}
  \item \emph{Say it:} \emph{kāpalādāru}, once more
  \item \emph{Say it:} say \emph{kāpalādāru}
  \item \emph{Recall:} say the Telugu for a plumber, then the Telugu for a mechanic, then say \emph{kāpalādāru} again
\end{itemize}

\begin{tcolorbox}[breakable,colback=teal!4,colframe=teal!35!black,title={Before you move on}]
What does \te{కాపలాదారు} mean? (\textquotedblleft{}A watchman, a guard\textquotedblright{}.) Say it once more.
\end{tcolorbox}
\section[\texorpdfstring{\te{పాలవాడు} (pālavāḍu)}{పాలవాడు (pālavāḍu)}]{\te{పాలవాడు} (pālavāḍu) — a milkman}

\label{lesson:TE-C314-palavadu}



\begin{quote}
Before the new one: say the Telugu for a mechanic, then the Telugu for a watchman.
\end{quote}

\subsection*{You'll want to know: \te{పాలవాడు}}
\textbf{\te{పాలవాడు}} — \emph{pālavāḍu} — \textquotedblleft{}a milkman\textquotedblright{}.

\begin{grammarlens}[title={What you've built}]
One more word for this chapter.
\end{grammarlens}

\subsection*{Your turn}
\begin{itemize}
\raggedright
  \item \emph{Say it:} \emph{pālavāḍu}
  \item \emph{Say it:} \emph{pālavāḍu}, once more
  \item \emph{Say it:} say \emph{pālavāḍu}
  \item \emph{Recall:} say the Telugu for a mechanic, then the Telugu for a watchman, then say \emph{pālavāḍu} again
\end{itemize}

\begin{tcolorbox}[breakable,colback=teal!4,colframe=teal!35!black,title={Before you move on}]
What does \te{పాలవాడు} mean? (\textquotedblleft{}A milkman\textquotedblright{}.) Say it once more.
\end{tcolorbox}
\section[(dialogue)]{Words from chapters 310-311 — a first pass}

\label{lesson:TE-R314-first-pass-words-from-chapters-310-311}



\begin{quote}
Say the words for a watchman and a milkman.
\end{quote}

\subsection*{Your turn}
Say these aloud:

\begin{itemize}
\raggedright
  \item a purse, a list, half, daily wages, an employee
  \item a college, a university, a degree, experience, an account
\end{itemize}

\begin{tcolorbox}[breakable,colback=teal!4,colframe=teal!35!black,title={Before you move on}]
A purse? (\textbf{\te{పర్సు}}, \emph{parsu}.) A list? (\textbf{\te{జాబితా}}, \emph{jābitā}.) Half? (\textbf{\te{సగం}}, \emph{sagaṁ}.) Daily wages? (\textbf{\te{కూలి}}, \emph{kūli}.) An employee? (\textbf{\te{ఉద్యోగి}}, \emph{udyōgi}.) A college? (\textbf{\te{కళాశాల}}, \emph{kaḷāśāla}.) A university? (\textbf{\te{విశ్వవిద్యాలయం}}, \emph{viśvavidyālayaṁ}.) A degree? (\textbf{\te{డిగ్రీ}}, \emph{ḍigrī}.) Experience? (\textbf{\te{అనుభవం}}, \emph{anubhavaṁ}.) An account? (\textbf{\te{ఖాతా}}, \emph{khātā}.)
\end{tcolorbox}
\section[(dialogue)]{Words from chapters 312-314 — a second pass}

\label{lesson:TE-R314-second-pass-words-from-chapters-312-314}



\begin{quote}
Say the words for a story and a game.
\end{quote}

\subsection*{Your turn}
Say these aloud:

\begin{itemize}
\raggedright
  \item a story, a game, a film, a coach, an engineer
  \item a piece of chalk, a slate, a slate pencil, ink, attendance
  \item a mason, a plumber, a mechanic, a watchman, a milkman
\end{itemize}

\begin{tcolorbox}[breakable,colback=teal!4,colframe=teal!35!black,title={Before you move on}]
A story? (\textbf{\te{కథ}}, \emph{katha}.) A game? (\textbf{\te{ఆట}}, \emph{āṭa}.) A film? (\textbf{\te{సినిమా}}, \emph{sinimā}.) A coach? (\textbf{\te{శిక్షకుడు}}, \emph{śikṣakuḍu}.) An engineer? (\textbf{\te{ఇంజనీరు}}, \emph{iñjanīru}.) A piece of chalk? (\textbf{\te{సుద్దముక్క}}, \emph{suddamukka}.) A slate? (\textbf{\te{పలక}}, \emph{palaka}.) A slate pencil? (\textbf{\te{బలపం}}, \emph{balapaṁ}.) Ink? (\textbf{\te{సిరా}}, \emph{sirā}.) Attendance? (\textbf{\te{హాజరు}}, \emph{hājaru}.) A mason? (\textbf{\te{మేస్త్రి}}, \emph{mēstri}.) A plumber? (\textbf{\te{ప్లంబరు}}, \emph{plambaru}.) A mechanic? (\textbf{\te{మెకానిక్కు}}, \emph{mekānikku}.) A watchman? (\textbf{\te{కాపలాదారు}}, \emph{kāpalādāru}.) A milkman? (\textbf{\te{పాలవాడు}}, \emph{pālavāḍu}.)
\end{tcolorbox}
